
If the Core Domain is a pristine, isolated mathematical cleanroom, the Infrastructure layer is the messy, volatile factory floor. This directory resides in the outermost ring of our Hexagonal Architecture. It is the only place in the codebase where side effects---such as reading from a hard drive, opening a TCP network socket, or allocating GPU memory---are legally permitted to occur.

The primary responsibility of the files in this directory is \textbf{adaptation}. They must import the pure Protocols defined in \texttt{core\_interfaces/} and implement the concrete, physical logic required to fulfill those contracts. Because they interact with the physical world, these files depend heavily on third-party libraries (\texttt{SQLAlchemy}, \texttt{requests}, \texttt{celery}, \texttt{transformers}) that are strictly banned from the inner rings.

To properly understand how these adapters function, we must first establish the mechanics of the third-party dependencies that power them.

% =====================================================================
% Core Infrastructure Dependencies
% =====================================================================

\subsubsection*{1. Object-Relational Mapping via \texttt{SQLAlchemy}}

At the core of our persistence infrastructure lies \texttt{SQLAlchemy}, the industry-standard SQL toolkit for Python. To understand its necessity, we must first examine a fundamental computer science problem known as the \textbf{Object-Relational Impedance Mismatch}.

Python and PostgreSQL operate on two entirely distinct mathematical paradigms for managing state. Python is an Object-Oriented language; it manages state dynamically in the heap memory as a graph of interconnected objects linked by memory pointers. PostgreSQL is a Relational Database Management System (RDBMS); it manages state on a hard drive using the strict rules of relational algebra, representing data as two-dimensional tabular sets (tables) linked by foreign keys.

Because these two paradigms are structurally incompatible, data cannot simply be copied from Python RAM to the PostgreSQL disk. It must be completely disassembled, translated into a string-based SQL query (e.g., \texttt{INSERT INTO ...}), transmitted over a TCP network socket, and then mathematically reconstructed by the database engine. When retrieving data, the exact reverse must occur.

An \textbf{Object-Relational Mapper (ORM)} is a software design pattern that entirely automates this translation layer. \texttt{SQLAlchemy}'s ORM allows developers to interact with the database using native Python classes and objects, rather than writing raw SQL strings. When we command the ORM to save a \texttt{DocumentMetadata} object, it dynamically parses the object's attributes, constructs the optimal SQL \texttt{INSERT} statement, and executes the network transaction.

Crucially, in our Hexagonal Architecture, the ORM is strictly an \textit{infrastructure detail}. The Core Domain is completely unaware that \texttt{SQLAlchemy} exists. The pure mathematical entities never inherit from ORM base classes, guaranteeing that our machine learning logic remains entirely decoupled from the database vendor.

To physically bridge this impedance mismatch, modern \texttt{SQLAlchemy} (version 2.0+) utilizes the \textbf{Declarative Mapping} paradigm. This allows us to define the database schema and the Python object behavior simultaneously within a single class definition. The core mechanics rely on a specific set of tools:

\begin{itemize}
    \item \textbf{\texttt{DeclarativeBase} (The Registry):} Before defining any tables, we create a master \texttt{Base} class that inherits from \texttt{DeclarativeBase}. This acts as a centralized metaclass registry. Every subsequent model we write will inherit from this \texttt{Base}. When the application starts, \texttt{SQLAlchemy} inspects this registry to know exactly which tables exist and how they mathematically relate to one another.
    \item \textbf{\texttt{Mapped[T]}:} This is a generic Python type hint. By declaring an attribute as \texttt{Mapped[str]}, we simultaneously satisfy Python's static type checkers (like \texttt{mypy}) and signal to the ORM that this attribute corresponds to a database column returning a string.
    \item \textbf{\texttt{mapped\_column()}:} While \texttt{Mapped} defines the Python-side behavior, \texttt{mapped\_column()} defines the \textit{database-side} constraints. We use it to specify configurations that have no pure Python equivalent, such as \texttt{primary\_key=True}, \texttt{nullable=False}, or declaring foreign key relationships.
    \item \textbf{Physical SQL Types (\texttt{String}, \texttt{Text}, \texttt{DateTime}):} Python simply sees a \texttt{str} or a \texttt{datetime} object. The database, however, requires strict memory allocation rules. We pass these explicit types into \texttt{mapped\_column()} to dictate the physical disk layout. For example, \texttt{String(255)} compiles to a bounded \texttt{VARCHAR(255)} (used for short titles), whereas \texttt{Text} compiles to an unbounded character object (required for long document abstracts).
\end{itemize}
By combining these primitives, a developer can define a complete PostgreSQL table simply by writing a type-hinted Python class.


Beyond defining the static schema, \texttt{SQLAlchemy} provides the active mechanics to query and mutate the database. The foundational engine for this is the \textbf{\texttt{Session}}. The \texttt{Session} acts as a temporary, localized workspace (or staging area) in Python RAM. When an application calls \texttt{session.add(document)}, the data is not immediately written to the disk. The \texttt{Session} tracks the object's state changes and handles the underlying TCP connection pooling, ultimately flushing the optimized SQL commands to the server only when instructed by the Unit of Work.

\textbf{The Composable Query Builder} \\
In standard SQL, extracting data requires writing a declarative string specifying what to get (\texttt{SELECT}), the conditions (\texttt{WHERE}), the sorting (\texttt{ORDER BY}), and the maximum number of results (\texttt{LIMIT}). In modern \texttt{SQLAlchemy} 2.0, querying is instead executed via a composable, object-oriented query builder.

By passing our mapped class into the \texttt{select()} function, we create a generative query object. We can then progressively apply filters and modifiers using Python method chaining:

\begin{lstlisting}[caption={Composable SQL in Python}]
stmt = (
    select(DocumentModel)
    .where(DocumentModel.title.ilike("%quantum%"))
    .order_by(DocumentModel.created_at.desc())
    .limit(10)
)
results = session.scalars(stmt).all()
\end{lstlisting}

This approach maps directly to standard relational algebra while maintaining strict Python type safety:
\begin{itemize}
    \item \textbf{\texttt{where()}:} Translates Python boolean logic into SQL \texttt{WHERE} clauses. Because \texttt{SQLAlchemy} overloads Python's comparison operators, an expression like \texttt{DocumentModel.id == X} does not evaluate to a standard \texttt{True} or \texttt{False} in Python; instead, it generates the Abstract Syntax Tree (AST) for the SQL equality check.
    \item \textbf{\texttt{order\_by()}:} Dictates the sorting mechanism. Attributes can be explicitly modified with \texttt{.desc()} (descending) or \texttt{.asc()} (ascending).
    \item \textbf{\texttt{limit()}:} Binds the result set to a maximum number of rows. This is structurally critical to prevent Out-Of-Memory (OOM) crashes when querying millions of documents.
    \item \textbf{\texttt{session.scalars(stmt)}:} The final execution step. It transmits the compiled, SQL-injection-proof string to PostgreSQL, receives the raw binary rows, and mathematically reconstructs them back into pure Python \texttt{DocumentModel} objects.
\end{itemize}

\textbf{PostgreSQL Dialects and Vendor-Specific Types} \\
While SQL is an international standard, every database vendor (e.g., PostgreSQL, MySQL, SQLite) implements custom, highly optimized features that exist outside the standard specification. To utilize these hardware-level optimizations, \texttt{SQLAlchemy} uses \textbf{Dialects}---translation engines that teach the ORM how to speak a specific database's exact flavor of SQL.

Because our infrastructure strictly relies on PostgreSQL, our Repository adapter explicitly imports PostgreSQL-specific data types to maximize memory efficiency and search speed:

\begin{itemize}
    \item \textbf{\texttt{PG\_UUID}:} (Imported via \texttt{from sqlalchemy.dialects.postgresql import UUID as PG\_UUID}). A naive relational database might store a UUID as a 36-character string (\texttt{VARCHAR}). PostgreSQL, however, has a native 128-bit \texttt{uuid} data type that dramatically reduces disk footprint and accelerates primary key index traversals. By using the dialect-specific \texttt{PG\_UUID}, we force the database to use this highly optimized binary format.
    \item \textbf{\texttt{ARRAY}:} Traditional relational algebra explicitly forbids storing lists of items inside a single column (violating First Normal Form). However, PostgreSQL provides a native \texttt{ARRAY} type. However, we will ultimately use the \texttt{pgvector} extension to execute the heavy vector calculus. %, understanding the \texttt{ARRAY} dialect is critical: it proves that PostgreSQL can natively understand and persist dense lists of primitives (like floating-point numbers) without requiring the engine to perform expensive, millions-of-rows table joins.
\end{itemize}


\subsubsection*{2. Database Migrations via \texttt{Alembic}}

While \texttt{SQLAlchemy} elegantly maps Python objects to database tables, it presents a temporal challenge. Software is not static. Over the lifecycle of the engine, the Core Domain will inevitably evolve; we might need to add a new metadata field, change a string length constraint, or attach an entirely new relational table.

If a developer simply modifies the \texttt{SQLAlchemy} \texttt{DeclarativeBase} class, the Python application will immediately expect the new schema, but the physical PostgreSQL database running on the disk will still possess the old schema. This structural mismatch causes fatal runtime exceptions. Furthermore, we cannot simply drop the tables and recreate them from scratch, as this would permanently destroy millions of processed documents.

To solve this, our infrastructure utilizes \textbf{\texttt{Alembic}}, a database migration tool built natively on top of \texttt{SQLAlchemy}. \texttt{Alembic} acts as a strict version control system for the physical database schema, functioning analogously to how Git tracks source code.

Instead of administrators executing raw SQL mutations on the production server, \texttt{Alembic} captures schema evolution as a Directed Acyclic Graph of migration scripts. Each script represents a deterministic state transition containing two functions:
\begin{itemize}
    \item \textbf{\texttt{upgrade()}:} Applies the structural change (e.g., executing an \texttt{ALTER TABLE ADD COLUMN} command).
    \item \textbf{\texttt{downgrade()}:} Reverts the exact change (e.g., \texttt{DROP COLUMN}), allowing the system to roll back seamlessly if a deployment fails.
\end{itemize}

Crucially, \texttt{Alembic} automates this via its \textit{autogenerate} engine. When triggered, \texttt{Alembic} inspects the \texttt{DeclarativeBase} metadata in Python RAM, connects to the live PostgreSQL database, and computes the topological difference between the two schemas. It then automatically authors the Python migration script required to synchronize the physical database with the codebase. By committing these generated scripts to our Git repository, we mathematically guarantee that any compute node pulling the code can deterministically reconstruct the exact physical database structure required to run the pipeline.


\textbf{The Physical Architecture of Alembic} \\
To achieve this version control, \texttt{Alembic} relies on a strict directory structure that dictates how the Python runtime interacts with the database:
\begin{itemize}
    \item \textbf{\texttt{alembic.ini}:} The global configuration file residing at the root of the repository. It defines the target database URL (\texttt{sqlalchemy.url}) and logging parameters.
    \item \textbf{\texttt{migrations/env.py}:} The critical execution context. This script boots up the \texttt{SQLAlchemy} engine and imports our \texttt{DeclarativeBase} metadata. It is the physical bridge that allows \texttt{Alembic} to read the Python models in RAM and compare them against the PostgreSQL tables on disk.
    \item \textbf{\texttt{migrations/script.py.mako}:} The template engine file. When \texttt{Alembic} autogenerates a new script, it uses this \texttt{Mako} template to structure the output file.
    \item \textbf{\texttt{migrations/versions/}:} The directory housing the actual Directed Acyclic Graph of historical state transitions.
\end{itemize}

\textbf{Handling Custom C-Extensions (The \texttt{pgvector} Challenge)} \\
When mapping standard types (\texttt{String}, \texttt{Integer}), \texttt{Alembic} automatically generates the correct \texttt{import sqlalchemy as sa} statements. However, when we introduce third-party C-extensions like \texttt{pgvector}, the autogenerator will blindly output a \texttt{Vector} column without importing the underlying library. If applied, the migration will crash with a \texttt{NameError}.

To mathematically prove the migration is sound, we must inject the custom dependency into the environment. We modify the \texttt{script.py.mako} template to automatically prepend \texttt{import pgvector.sqlalchemy} to the top of every new migration file. Furthermore, we must ensure our \texttt{env.py} file correctly imports our \texttt{models.py} file so the \texttt{target\_metadata} is perfectly synced with our domain logic.

\textbf{The Migration Lifecycle in Practice} \\
Consider a scenario where we expand our \texttt{DocumentModel} to include a 768-dimensional tensor for machine learning embeddings. We update our Python model with \texttt{mapped\_column(Vector(768))}.

We then instruct \texttt{Alembic} to calculate the topological difference and generate the script via the terminal:

\begin{lstlisting}[language=bash, caption={Generating the Migration}]
$ alembic revision --autogenerate -m "feat_add_vec_embedding_column"
\end{lstlisting}

This generates a new file in the \texttt{versions/} directory with a unique cryptographic hash (e.g., \texttt{21178132d92e\_...py}). The autogenerated script clearly defines the bidirectional state transition:

\begin{lstlisting}[caption={The Deterministic State Transition (Simplified)}]
from alembic import op
import sqlalchemy as sa
import pgvector.sqlalchemy # Injected custom type

# Revision identifiers
revision = '21178132d92e'
down_revision = '51ab2670c483' # Pointer to the previous state

def upgrade() -> None:
    # Move state forward: Add the vector column
    op.add_column(
        'documents',
        sa.Column('embedding', pgvector.sqlalchemy.Vector(dim=768))
    )

def downgrade() -> None:
    # Revert state backward: Destroy the vector column
    op.drop_column('documents', 'embedding')
\end{lstlisting}

Notice the \texttt{down\_revision} variable. This establishes the strict mathematical edges of our Directed Acyclic Graph; this script can \textit{only} be applied if the database is currently sitting at state \texttt{51ab2670c483}.

Finally, to physically mutate the PostgreSQL schema, we execute the upgrade command, pushing the database to the most recent node (\texttt{head}) in the graph:

\begin{lstlisting}[language=bash, caption={Applying the State Transition}]
$ alembic upgrade head
\end{lstlisting}

\textbf{Connection Strings and Environment Injection} \\
To physically execute these migrations, \texttt{Alembic} (and \texttt{SQLAlchemy}) must establish a TCP connection to the PostgreSQL server. This is defined by a Data Source Name (DSN) URL. Consider the default string found in our \texttt{alembic.ini}:

\begin{lstlisting}[caption={The Database Connection String}]
sqlalchemy.url = postgresql+psycopg://semantic_user:
                 password@localhost:5432/semantic_engine_db
\end{lstlisting}

This string mathematically dictates the exact network path and protocol stack:
\begin{itemize}
    \item \textbf{\texttt{postgresql+psycopg}:} The Dialect and the DBAPI. It tells the engine to speak PostgreSQL using the \texttt{psycopg} (version 3) C-driver.
    \item \textbf{\texttt{semantic\_user:password}:} The authentication credentials.
    \item \textbf{\texttt{localhost:5432}:} The network location (host and port) of the database server.
    \item \textbf{\texttt{semantic\_engine\_db}:} The specific database instance to connect to.
\end{itemize}

However, a critical architectural conflict arises here. The \texttt{alembic.ini} file is committed to version control. If we hardcode our production database credentials into this file, we create a severe security vulnerability. Furthermore, a hardcoded URL destroys system portability; the CI/CD hypervisor running tests on GitHub Actions needs to connect to a temporary test database, while a developer's local machine needs to connect to a local Docker container.

To resolve this, we utilize \textbf{Environment Variable Injection} within \texttt{migrations/env.py}. We intercept \texttt{Alembic}'s boot sequence and overwrite the static configuration dynamically:

\begin{lstlisting}[caption={Dynamic URL Injection in \texttt{env.py}}]
import os
from logging.config import fileConfig
from sqlalchemy import engine_from_config
from alembic import context

config = context.config

# 1. Read the dynamic URL from the host operating system
database_url = os.getenv("DATABASE_URL")

# 2. Overwrite the static alembic.ini value in memory
if database_url:
    config.set_main_option("sqlalchemy.url", database_url)
\end{lstlisting}

This pattern guarantees that the codebase remains mathematically invariant across all environments. The URL in \texttt{alembic.ini} serves strictly as a local fallback for development. When deployed to Production, the orchestrator (e.g., Kubernetes or Docker Compose) securely passes the production \texttt{DATABASE\_URL} into the container's RAM. \texttt{env.py} reads this ephemeral environment variable and transparently redirects \texttt{Alembic} to the production server without altering a single line of committed code.


\subsubsection*{3. Vector Mathematics via \texttt{pgvector}}


% If we attempted to find the semantic distance between documents using pure Python, the system would have to execute a \texttt{SELECT *}, load $1$ million vectors from the disk into the server's RAM, and then compute the dot products sequentially. This would immediately trigger a fatal Out-Of-Memory (OOM) exception.


Standard relational databases are highly optimized for B-Tree indexing and boolean scalar operations. They are fundamentally not designed to perform linear algebra on dense, 768-dimensional floating-point tensors.

\textbf{The CPU vs. GPU Compute Boundary} \\
In modern machine learning pipelines, there is a strict hardware division of labor. Generating a 768-dimensional embedding from raw text requires passing tokens through billions of neural network parameters; this highly parallelized matrix multiplication is executed on a \textbf{GPU} using frameworks like \texttt{PyTorch}. We will define this process later in the ML Adapter layer.

However, once those embeddings are generated and stored, \textit{searching} them is a completely different computational problem. If a user queries the engine, we must calculate the mathematical distance between the user's query vector and $1$ million document vectors stored in the database. If we attempted to send these $1$ million vectors from the database RAM across the motherboard's PCIe bus to the GPU for calculation, the I/O transfer time would completely eclipse the computation time, creating a massive bottleneck.

% To solve this, we push the search computation directly down to the storage layer using \textbf{\texttt{pgvector}}, an open-source C-extension for PostgreSQL. Instead of moving the data to the compute (GPU), we move the compute to the data (CPU). \texttt{pgvector} physically alters the database engine, allowing it to store tensors natively and execute highly optimized distance calculations directly on the database server's CPU using SIMD (Single Instruction, Multiple Data) hardware instructions like AVX-512. The Python application only ever receives the final, lightweight list of the top 10 closest document IDs.

To solve this, we push the mathematical computation down to the storage layer using \textbf{\texttt{pgvector}}, an open-source C-extension for PostgreSQL. \texttt{pgvector} physically alters the database engine, allowing it to store tensors natively and execute highly parallelized distance calculations (Cosine distance, Euclidean $L^2$ distance) directly on the disk using CPU SIMD (Single Instruction, Multiple Data) instructions. The Python application only ever receives the final, lightweight list of the top 10 closest document IDs.

To utilize this in our ORM, we import the custom type definition:

\begin{lstlisting}[caption={Mapping the Vector Type}]
from pgvector.sqlalchemy import Vector # type: ignore
# ... other ...
class DocumentModel(Base):
    __tablename__ = "documents"
    # ... other columns ...
    embedding: Mapped[list[float]] = mapped_column(Vector(768))
\end{lstlisting}

\textbf{Type Theory: Bypassing the Static Analyzer} \\
Notice the explicit \texttt{\# type: ignore} directive appended to the import statement. In our Core Domain, we relied heavily on static type checkers (like \texttt{mypy}) to provide mathematical proofs of our contracts. However, static analyzers operate entirely by reading Python source code (\texttt{.py} files or \texttt{.pyi} interface stubs) at compile-time.

Because \texttt{pgvector} is a C-extension, its core logic is compiled directly into a binary shared object file %(\texttt{.so} or \texttt{.dll})
that is dynamically injected into the interpreter at runtime. The static analyzer cannot inspect compiled C code. Consequently, it will flag the import as a fatal \texttt{ModuleNotFoundError} or an \texttt{Any} type violation. By explicitly injecting the \texttt{ignore} directive, we are deliberately asserting a boundary condition: we are informing the type checker that we mathematically guarantee the existence of this binary extension at runtime, overriding the static proof system for this specific hardware optimization.
% TODO: Explain why we use `# type: ignore` in Python when dealing with dynamically injected C-extensions.

\subsubsection*{4. Cryptographic Identity via \texttt{UUID}}
% TODO: Explain the difference between the `uuid` module and `uuid4`. Why pure randomness (v4) is critical for distributed systems to avoid primary key collisions without locking the database.

% =====================================================================
% The Concrete Adapters
% =====================================================================

\subsubsection*{5. The Database Models (\texttt{database/models.py})}
% TODO: Map the pure Domain Entity to SQLAlchemy Declarative Base.

\subsubsection*{6. The Repository Adapter (\texttt{database/repository.py})}
% TODO: Implement the DocumentRepository Port using the SQLAlchemy ORM.

\subsubsection*{7. The Unit of Work Adapter (\texttt{database/uow.py})}
% TODO: Implement the AbstractUnitOfWork Port using SQLAlchemy sessionmaker and actual commits/rollbacks.

\subsubsection*{8. The API Adapter (\texttt{api/arxiv.py})}
% TODO: Implement the AcademicGraphPort using HTTP requests and XML parsing to fetch arXiv papers.

\subsubsection*{9. The ML Adapter (\texttt{ml/sentence\_transformer.py})}
% TODO: Implement the TextEmbeddingPort using HuggingFace models to generate the 768-dimensional vectors.

\subsubsection*{10. The Distributed Workers (\texttt{worker/celery\_app.py} and \texttt{tasks.py})}
% TODO: Explain how Celery acts as the asynchronous event loop, pulling jobs from Redis to execute the adapters across multiple hardware nodes.
